\scanpage{340}
\mspage{39}

aucune exception. Instaurée par Dieu dès avant la Création des âmes et celle de l'Univers, elle s'enracine dans la nature même de Dieu et de Sa relation à Lui-même.

Il est à noter ici que même dans le cas d'une fille, la mère est bel et bien investie du rôle d'Epouse à côté de celui de la Mère. Et de même dans le cas d'un garçon, le père est, en plus de son rôle de Père, investi de celui d'Epoux. Il faut interpréter ces faits fondamentaux non comme le signe d'une soi-disante nature ``asexuée'', ou sexuellement indifférenciée, du petit enfant, mais au contraire comme attestant de la nature foncièrement ``\textit{bisexuée}'' ou ``\textit{androgyne}'' de la psyché. Ce caractère essentiel de l'âme agit tout au cours de la vie depuis le bas âge jusque dans la vieillesse, et au delà même de la mort, dans l'au-delà.

Il est vrai que l'âme, tant dans sa pulsion amoureuse que par ses qualités maîtresses, a été créée soit féminine, soit masculine, et qu'au cours d'une existence terrestre, cette qualité première se reflète le plus souvent (mais nullement toujours) dans le sexe biologique. Cela n'empêche que dans tous les cas sans exception, l'âme participe également aux forces et aux qualités du sexe opposé. Aussi la pulsion du sexe l'attire-t-elle non seulement vers les êtres de l'autre sexe, mais aussi vers ses congénères. Cette bisexualité foncière de l'âme est sanctionnée par cette autre loi spirituelle, que toute âme de femme est appelée à s'incarner non seulement comme femme, mais aussi de temps à autre comme homme. Et inversement, toute âme d'homme s'incarne parfois en des existences de femme.

Cette qualité bisexuée de l'âme est exprimée également par son ``\textit{corps éternel}'', ou plus précisément par le fait qu'en même temps qu'Il créa l'âme, Dieu a créé aussi pour elle \textit{deux} ``corps éternels'' : l'un corps de femme, l'autre corps d'homme, et l'un et l'autre sa propriété inaliénable en toute éternité. Ce sont des corps parfaits, immatériels, les ``\textit{corps de lumière}'' dont parle la tradition mystique. Au cours du temps d'errance de l'âme, leur rôle est de servir de ``maquette'' aux corps qui seront les siens dans ses différentes incarnations soit comme femme, soit comme homme. Il est d'ailleurs des états de la psyché où, par Acte de Dieu, il nous est donné d'avoir la perception de l'un ou l'autre de nos deux corps éternels, comme si nous l'habitions déjà bel et bien et que nos mouvements et notre vie au fil des instants s'accomplissent en lui. Ainsi, depuis la ``montée des forces féminines'' dans mon être, en juin 1976, j'ai la perception permanente de mon corps féminin, perception plus ou moins vive d'un moment à l'autre, ou d'une période à l'autre de ma vie.

Tout ceci nous montre qu'à côté de la pulsion amoureuse ``orthodoxe'', celle qui attire la fille ou la femme vers l'homme, et le garçon ou l'homme vers la femme (et souvent recouverte par elle, ou s'amalgamant à elle pour échapper au regard...), la \textit{pulsion amoureuse ``homosexuelle'' vit dans toute âme humaine, tout comme elle vit en Dieu.} Elle fait partie des forces créatrices profondes de l'âme, toute aussi essentielle que la pulsion amoureuse orthodoxe, et comme elle destinée à nourrir ses élans créateurs en toute éternité.

Sous l'effet de l'opprobre sociale qui, de nos jours encore, pèse sur elle, cette pulsion ``maudite'' est le plus souvent refoulée dans l'inconscient moyen et profond, par les soins des mécanismes discrets et efficaces de la \textit{répression sexuelle}, intériorisée par la psyché. Découvrir l'existence en elle de cette pulsion, à l'encontre des tabous puissants qui pèsent sur elle, et assumer sa présence et la beauté du désir qu'elle inspire, fait partie des tâches essentielles de l'âme pendant son temps d'errance.

% archive sheet number: 340
